\documentclass[11pt]{article}
\usepackage{amsmath,amssymb}
\usepackage[margin=1in]{geometry}
\usepackage{hyperref}
\emergencystretch=3em

\title{Quotients in distribution and the leading-order posterior expectation (Programme S, tranche A4-L)}
\author{Claude, with Daniel Murfet}
\date{2026-09-08}

\begin{document}
\maketitle
\sloppy

\begin{abstract}
Units 286--288 close Programme S with the leading term of \texttt{cor:empirical\_expectation} after
finite chart assembly. The generic statement is probabilistic: if $(Y_\ell,W_\ell)\Rightarrow(A,B)$
jointly and $P(B=0)=0$ then $Y_\ell/W_\ell\Rightarrow A/B$ (Lean's totalised division; nothing about
$W_\ell\ne0$ at finite $\ell$). Mathlib has no continuous mapping theorem for maps continuous only
almost surely, so the proof uses the truncated division $T_\delta(a,b)=ab/\max(b^2,\delta^2)$, globally
continuous and equal to $a/b$ on $|b|\ge\delta$: for a bounded Lipschitz test function of oscillation
$C$ the integrals of $f(Y/W)$ and $f(T_\delta(Y,W))$ differ by at most $C\,P(|W|\le\delta)$, portmanteau
bounds $\limsup P(|W_\ell|\le\delta)$ by $P(|B|\le\delta)$, and $P(|B|\le1/(m+1))\downarrow P(B=0)=0$.
The assembled corollary takes the pair (denominator, numerator) of joint chart data, a.e.\ chart
decompositions with residuals negligible at the leading scale $a_\ell=N_\ell^{-\mu_0}(\log N_\ell)^{j_0}$,
aggregate predecessor sums vanishing a.s.\ at every source index, and $P(C^1_{\mu_0,j_0}(X)=0)=0$; the
normalised pair converges jointly by vector Slutsky and the quotient theorem finishes (Headline~XXXVIII).
Review v31: pass. Three units against a cap of twelve. No \texttt{sorry}, no additional \texttt{axiom};
3415 jobs; pin \texttt{a1a7d1a}.
\end{abstract}

\section{Why only the leading term}
Astra~\#35 refused the all-orders quotient expansion: \texttt{lemma:division} has no statement in the
source, and dividing by $1+c/\log n$ produces an infinite block of inverse logarithmic powers at a fixed
power of $n$, which the finite predecessor ordering of the Taylor tree cannot host. The leading term,
however, is exactly a joint-convergence-plus-division statement, and the assembled Taylor tree already
provides the joint convergence.

\section{The generic quotient theorem (u286)}
The weak convergence of laws is characterised by integrals of bounded Lipschitz test functions
(\texttt{tendsto\_iff\_forall\_lipschitz\_integral\_tendsto}, hence the countably generated filter). For
such an $f$ with oscillation bound $C$: (i) $|f(Y/W)-f(T_\delta(Y,W))|\le C\,\mathbf 1_{\{|W|\le\delta\}}$
pointwise, so the integral difference is at most $C\,P(|W|\le\delta)$ (indicator integral on a probability
space; integrability from the oscillation bound, $\|f(x)\|\le\|f(0)\|+C$); (ii) $W_\ell\Rightarrow B$ by
projecting the pair, and \texttt{ProbabilityMeasure.limsup\_measure\_closed\_le\_of\_tendsto} on the closed
set $\{|b|\le\delta\}$ gives eventually $P(|W_\ell|\le\delta)<\eta$ once $P(|B|\le\delta)<\eta$; (iii) the
sets $\{|B|\le1/(m+1)\}$ decrease to $\{B=0\}$, so $\delta$ exists (\texttt{tendsto\_measure\_iInter\_atTop});
(iv) continuous mapping at fixed $\delta$; (v) $\varepsilon/3$ with $\eta=\varepsilon/(6(C+1))$. The scaled
form removes a common deterministic normalisation $a_\ell\ne0$ (the cancellation
$(U/a)/(V/a)=U/V$ fails at $a=0$ under totalised division, so $a_\ell\ne0$ is required at every index).

\section{The assembled leading quotient (u287--288)}
\texttt{PairData} is a \texttt{def} synonym of \texttt{JointData × JointData} with the sup norm and the
Borel $\sigma$-algebra (denominator first, numerator second). The type does not enforce a common phase
or $\eta_\phi=\phi\cdot\eta_1$; the paper's model is a specialisation. With the a.e.\ decompositions and
vanishing predecessor sums,
$Z^0_\ell/a_\ell-C(X_\ell)=(R_\ell-C(X_\ell))+E_\ell/a_\ell$ a.e.\ for both families; the first term is
uniformly small on joint balls (the sup norm of the pair dominates both components), the second is the
normalised residual; both pairs tend to zero in probability (two small lemmas: sums and pairs of
in-probability-zero sequences), and vector Slutsky gives $(Z^0[\phi]/a,Z^0[1]/a)\Rightarrow(C^\phi(Y),C^1(X))$.
The reviewer's points: remainders minus coefficients \emph{at the same source datum} tend to zero (the
limit lives on another probability space); aggregate predecessor sums vanish, individual coefficients
need not; the numerator coefficient may vanish; $P(C^1=0)=0$ is not positivity; no convergence of
expectations follows.

\section{Lessons}
\begin{itemize}
\item \textbf{Truncate the map, not the variables.} A globally continuous surrogate agreeing with the
discontinuous map off a small-probability set replaces the a.s.-continuous mapping theorem; the test
function's boundedness pays for the bad set.
\item \textbf{\texttt{rw} versus \texttt{refine} on weak-convergence characterisations.} Rewriting the
\texttt{Tendsto} goal with the characterisation fails on instance mismatch; \texttt{refine (lemma (Ω := ℝ)).2 ?\_}
unifies.
\item \textbf{Order of the pair.} Store data as (denominator, numerator) and produce coefficients as
(numerator, denominator) for division; the reviewer checked this explicitly --- document it.
\item \textbf{Product $\sigma$-algebras.} Products of nonseparable function spaces need a declared Borel
instance on a synonym; never let \texttt{Pi.borelSpace}/\texttt{Prod.borelSpace} be assumed.
\end{itemize}
\end{document}
